% Short document: no title page. The title block opens page one and the
% document goes straight into its prose. Use this for anything up to about
% twelve pages; past that, blank-template.tex (a dedicated title page and a
% contents page). Build: lualatex, three passes.
% Canonical preamble for this skill. See references/latex-house-style.md.
% The look is references/house-style/ (style-spec.md, housestyle.sty); this
% file adds only the teaching affordances and the math kit on top of it.
% Replace <COURSE>, <DOC>, <NAME> and <YEAR> as needed.
% It ends before \begin{document}, so a smoke document can \input it as-is.
% Build with scripts/build.sh (lualatex, three passes); it puts housestyle.sty
% on the search path and the .sty finds the vendored fonts itself.

\documentclass[11pt]{article}
% A4: uncomment the next line; the margins stay the spec's.
% \PassOptionsToPackage{a4paper}{geometry}
\usepackage{housestyle}  % geometry, fonts, colour tokens, running head, the twelve blocks, hyperref (hidelinks)
\usepackage{amssymb,amsthm}
\usepackage{mathtools}
\usepackage{tabularx,multirow}

% --- Running head and foot -------------------------------------------------
% Slug at the left of the head; the section name at the right follows each
% \section (override with \hssection{...}). Foot: the copyright line at the
% left, the page number at the right. The first page has no head, per the spec.
\hsslug{<COURSE> <DOC>}
\fancyfoot[L]{\hslabel{\textcopyright{} <YEAR> <NAME>. All rights reserved.}}
\makeatletter
\renewcommand{\sectionmark}[1]{\gdef\@hssection{#1}}
\makeatother

% --- Headings --------------------------------------------------------------
% Headings are unnumbered on the page (the .sty's \section). Cross-reference
% a section by name with \nameref{sec:...}; a Part by its number, Part~\ref{part:...}.
% A Part opens on a new page: its eyebrow, then the page heading at 28/34, no bold.
\makeatletter
\renewcommand{\thepart}{\arabic{part}}
\newcommand{\hs@parteyebrow}{}
\titleclass{\part}{top}
\assignpagestyle{\part}{fancy}  % keep the head and the copyright foot on a Part page
\titleformat{\part}[display]{\fontsize{28pt}{34pt}\selectfont\color{ink}}%
  {\hsaccentlabel{Part \thepart\hs@parteyebrow}}{10pt}{}
\titlespacing*{\part}{0pt}{0pt}{24pt}
\titleformat{\subsection}{\fontsize{14pt}{19pt}\selectfont\color{ink}}{}{0pt}{}
\titlespacing*{\subsection}{0pt}{18pt}{6pt}
\titleformat{\subsubsection}{\itshape\color{ink}}{}{0pt}{}
\titlespacing*{\subsubsection}{0pt}{12pt}{2pt}

% --- Taxonomy: category as structure, never as colour ------------------------
% \catbanner{A}{Title} opens a Part: eyebrow "Part N . Category A", then the
% heading, and the running head's section becomes the title.
\newcommand{\catbanner}[2]{%
  \gdef\hs@parteyebrow{\ \textperiodcentered\ Category #1}%
  \part{#2}\hssection{#2}\gdef\hs@parteyebrow{}}
\makeatother
\newcommand{\cat}[1]{\hslabel{[#1]}}      % [A]
\newcommand{\tool}[1]{\hslabel{[T#1]}}    % [T5]
\newcommand{\tools}[1]{\hslabel{[#1]}}    % [T4,T5]

% --- Teaching affordances, inside the twelve blocks -------------------------
% insight: the section's lead line, right under the heading.
\newcommand{\insight}[1]{\hslead{#1}}
% derivation: a display-equation run (block 12) the reader can skip on
% re-read. A label opens it, a grey rule shows where it ends.
\newenvironment{derivation}%
  {\par\vspace{4pt}\hslabel{Derivation \textperiodcentered\ skip on re-read}\par\nopagebreak}%
  {\par\nopagebreak\vspace{2pt}\noindent\hsthinrule\par}
% \onsheet: this equation lives on the formula sheet. Grey, not accent.
\newcommand{\onsheet}{\par\nopagebreak\noindent\hslabel{On the formula sheet}\nopagebreak}
% \opt: out of scope. Inline after a mention, or first thing in an optional paragraph.
\newcommand{\opt}{\leavevmode\unskip\hspace{4pt}\hslabel{[optional]}\hspace{4pt}\ignorespaces}
% \probhead{source}{method}{category}: one mono label line over a worked problem.
\newcommand{\probhead}[3]{%
  \par\vspace{8pt}\noindent\hslabel{#1 \textperiodcentered\ #2 \textperiodcentered\ #3}\par\nopagebreak}

% --- Tables ----------------------------------------------------------------
% L{w} and R{w} come from the .sty; Y is a ragged-right tabularx column.
\newcolumntype{Y}{>{\raggedright\arraybackslash}X}

% --- Math kit ----------------------------------------------------------------
\newcommand{\R}{\mathbb{R}}
\newcommand{\N}{\mathbb{N}}
\newcommand{\Z}{\mathbb{Z}}
\newcommand{\RR}[1]{\mathbb{R}^{#1}}
\newcommand{\grad}{\nabla}
\newcommand{\Hess}{H_f}

% Bold vectors -- define all lowercase so agents can grab any letter
\newcommand{\ba}{\mathbf{a}}  \newcommand{\bb}{\mathbf{b}}  \newcommand{\bc}{\mathbf{c}}
\newcommand{\bd}{\mathbf{d}}  \newcommand{\be}{\mathbf{e}}  \newcommand{\bff}{\mathbf{f}}
\newcommand{\bg}{\mathbf{g}}  \newcommand{\bh}{\mathbf{h}}  \newcommand{\bi}{\mathbf{i}}
\newcommand{\bj}{\mathbf{j}}  \newcommand{\bk}{\mathbf{k}}  \newcommand{\bl}{\mathbf{l}}
\newcommand{\bm}{\mathbf{m}}  \newcommand{\bn}{\mathbf{n}}  \newcommand{\bp}{\mathbf{p}}
\newcommand{\bq}{\mathbf{q}}  \newcommand{\br}{\mathbf{r}}  \newcommand{\bs}{\mathbf{s}}
\newcommand{\bt}{\mathbf{t}}  \newcommand{\bu}{\mathbf{u}}  \newcommand{\bv}{\mathbf{v}}
\newcommand{\bw}{\mathbf{w}}  \newcommand{\bx}{\mathbf{x}}  \newcommand{\by}{\mathbf{y}}
\newcommand{\bz}{\mathbf{z}}  \newcommand{\bzero}{\mathbf{0}}

% Transpose: braces (not ^) so X^\trans gives X^{\!\top}
\newcommand{\trans}{{\!\top}}

% Partial derivatives
\newcommand{\pderiv}[2]{\dfrac{\partial #1}{\partial #2}}
\newcommand{\ppderiv}[2]{\dfrac{\partial^2 #1}{\partial #2^2}}
\newcommand{\pmderiv}[3]{\dfrac{\partial^2 #1}{\partial #2 \, \partial #3}}

% Big-O -- use \Oh{n} or \Ohof{n \log n} everywhere; never bare $O()$
\newcommand{\Oh}{\mathcal{O}}
\newcommand{\Ohof}[1]{\mathcal{O}\!\left(#1\right)}

% Differential d
\newcommand{\diff}[1]{\mathop{}\!\mathrm{d}#1}

% --- Code --------------------------------------------------------------------
% Block 11: \hslisting{Listing 1 / what it shows} then
% \begin{Verbatim}[bgcolor=codefill] ... \end{Verbatim}. No syntax colour in print.

% \hsnumbersections  % uncomment for numbered sections (1, 1.1, ...)
\hsslug{<Slug>}
\begin{document}

\hstitleblock{<Eyebrow> \textperiodcentered\ <Month Year>}{<Document title>}%
{<A one-sentence lead that says what this document is for.>}

The opening paragraph is the only general one. It says what the thing is and
why a reader should care, before any number or mechanism appears.

A second paragraph can set up the figure below. Nothing is illustrated twice:
if the diagram carries it, the prose does not repeat it.

\begin{hsfigure}{Figure 1 / node-and-arrow flow}%
{Slate boxes are inputs and outputs. Sage boxes are work an agent does.
Rust boxes are deterministic checks that pass or fail.}
\begin{tikzpicture}[node distance=14pt]
  \node[hsnode] (a) {\hsnodetext{Input}{A brief}{goal, boundaries, done-criteria}};
  \node[hswork,right=of a] (b) {\hsnodetext{Agent}{Draft}{structure first, then prose}};
  \node[hsgate,right=of b] (c) {\hsnodetext{Gate}{Style check}{pass, or back with reasons}};
  \node[hsnode,right=of c] (d) {\hsnodetext{Output}{A PDF}{with its sources attached}};
  \draw[hsarrow] (a) -- (b); \draw[hsarrow] (b) -- (c); \draw[hsarrow] (c) -- (d);
  \draw[hsfail] (c.south) -- ++(0,-12pt) -| (b.south);
\end{tikzpicture}
\hsfailnote{A failed gate returns to the draft step, twice, then to a person.}
\end{hsfigure}

\hsprovenance{Figure 1 from <file>, read <date>.}
\newpage

\section{Measured results}
\hslead{Two to four measured figures, each saying what was measured and when.}

\begin{hsstatrow}[3]
\hsstat{42}{What was counted, and on which date.}
\hsstat{7.5}{Caption for the second figure.}
\hsstat{100\%}{Caption for the third figure.}
\end{hsstatrow}

\hsclaim{A single pull-quote claim, set large between two rules.}{Attribution label}

\begin{hscallout}{What this does not mean}
A callout holds the objection a careful reader would raise, stated before they
have to. Placeholder text.
\end{hscallout}

\hsprovenance{Figures read on <date> from <command or file>.}
\hssourcespage
\begin{hssources}
\item Author, \emph{Title of the first source}, Publisher, Year.
\item Author, \emph{Title of the second source}. \hsurl{https://example.com/source}
\end{hssources}
\end{document}
